\section{Fragestellung: was ist FFGFTs Zeit-Windung genau?}
In FFGFT ist die Zeit der Massenkreis, über $\tilde T\cdot m=1$ in die übrigen Kreise eingewoben
und kein abspaltbarer Faktor (Dok.~285). Daran schließt eine präzise Frage an: was genau
\emph{ist} FFGFTs Zeit-Windung?

Dieses Dokument klärt sie vollständig FFGFT-intern -- sie ist aus der fraktalen Korrektur
\emph{ableitbar}. Das Ergebnis in einem Satz: die K$_{\text{frak}}$-Rekursion zwingt die
Zeit-Windung zu einer \emph{logarithmischen Spirale mit Selbstähnlichkeits-Verhältnis $e$}, also
zu kontinuierlicher Skaleninvarianz.

\section{Drei Bausteine aus dem Corpus}
Drei Größen sind an anderer Stelle festgelegt und werden hier nur zusammengeführt. Erstens die
fraktale Dimension des Raums, $D_f=3-\xi_0$ mit $\xi_0=4/30000$ (Dok.~205, Dok.~101).
Zweitens die fraktale Korrektur im Gatter: die X-Drehung erfolgt nicht um $\pi$, sondern um
$\alpha=\pi\,K_{\text{frak}}$ mit $K_{\text{frak}}=1-100\,\xi_0$ (Dok.~034). Drittens der
Skalenlauf als Rekursion $\xi_{n+1}=\xi_n\,(1-100\,\xi_0)$, genauer
$\xi_{n+1}=\xi_n(1-100\,\xi_n)$ (Dok.~146).

Aus $\xi_0=1/7500$ folgen drei exakte Tatsachen, die im Weiteren tragen:
\begin{equation}
  100\,\xi_0=\frac{1}{75},\qquad K_{\text{frak}}=\frac{74}{75},\qquad
  \frac{1}{100\,\xi_0}=75 .
\end{equation}

\section{Nichtschließung: von der 75-Schließung zur Spirale}
Eine Windung, die pro Umlauf um $2\pi$ vorrücken sollte, rückt gatterweise nur um
$\pi\,K_{\text{frak}}$ vor und \emph{verfehlt} $\pi$ um $100\pi\,\xi_0$ je Halbdrehung.
Iteriert man, schließt die Windung nicht, sondern präzediert -- sie wird zur Spirale. Der
Schließungs-Defekt je Umlauf beträgt $d=100\,\xi_0$ in Umlauf-Einheiten.

Bei eingefrorenem $\xi_0$ ist $d=1/75$ konstant, die Rotationszahl $74/75$ also rational: die
Windung schließt nach \emph{genau $75$ Umläufen} und trägt eine $75$-zählige Struktur -- die
Zahl $75=1/(100\,\xi_0)$ fällt exakt heraus. Der Grenzfall konstanter Ganghöhe,
$d=0$, schließt dagegen bereits nach einem Umlauf. Beide
geschlossenen Fälle sind die \emph{aufgerollte}, kompakte Lesart der Zeit-Windung: eine einzige
Skala, die Windung liegt geschlossen und ist statisch. FFGFT
kennt eine Nichtschließung bereits algebraisch, in der nicht-multiplikativen Modenkomposition
(Dok.~193); hier tritt ihr geometrisch-temporales Gegenstück hinzu.

\section{Der Skalenlauf macht die Spirale logarithmisch}
Läuft $\xi$ über die Rekursion, so verwandelt sich die rationale $75$-Schließung in eine
nicht-schließende, selbstähnliche Spirale. Dieser Übergang ist die \emph{Entrollung} der Windung:
die Spirale ist das dekompaktifizierte Gegenstück der geschlossenen $75$-Windung, und
\enquote{entrollen} heißt hier präzise, den Skalenlauf laufen zu lassen (die Rekursion), nicht
bloß einen Torusfaktor abzuspalten. Die Ableitung ist geschlossen. Mit $u_n=1/\xi_n$
ist
\begin{equation}
  u_{n+1}=\frac{1}{\xi_n(1-100\,\xi_n)}=u_n\,(1+100\,\xi_n+\mathcal{O}(\xi_n^2))
         =u_n+100+\mathcal{O}(\xi_n),
\end{equation}
also $u_n\approx 100\,n+7500$ und damit ein \emph{algebraischer} $1/n$-Zerfall
\begin{equation}
  \xi_n\;\approx\;\frac{1}{100\,(n+75)} .
\end{equation}
Der kumulierte Präzessions-Defekt wird dadurch logarithmisch,
\begin{equation}
  D(k)=\sum_{n\le k}100\,\xi_n\;\approx\;\sum_{n\le k}\frac{1}{n+75}
       \;\approx\;\ln\!\Big(1+\frac{k}{75}\Big).
\end{equation}
Mit Radius $\rho=k+75$ -- der Pol der Spirale liegt beim effektiven Ursprung
$n=-75=-1/(100\,\xi_0)$, derselben $75$ wie zuvor -- und Präzessionswinkel $\beta=2\pi\,D(k)$
folgt $\rho/75=e^{\beta/2\pi}$, also
\begin{equation}
  \rho\;=\;75\,e^{\beta/2\pi}\;=\;\text{logarithmische (äquiangulare) Spirale.}
\end{equation}
Ihre Selbstähnlichkeit ist damit festgelegt: je Präzessionsrunde, $\Delta\beta=2\pi$, skaliert
der Radius um $e^{1}=e$. Das \emph{Verhältnis $e$ ist nicht angesetzt, sondern folgt} aus dem
algebraischen $1/n$-Zerfall der Rekursion: es ist die Basis des natürlichen Logarithmus, den die
harmonische Summe erzeugt, und kein zusätzlicher Parameter tritt hinzu. Auch der Radius-Offset
$75$ ist nicht frei gewählt -- er ist $1/(100\,\xi_0)$, dieselbe Konstante, die schon die
$75$-Schließung trägt.

Numerisch bestätigt sich beides ohne Spielraum (\texttt{zeit\_logspirale\_probe.py}): der Fit
$\ln\rho=a\,\beta+b$ gibt $a=0{,}15922\approx 1/2\pi=0{,}15915$ und $e^{2\pi a}=2{,}7195\approx e$;
das lokal gemessene Selbstähnlichkeits-Verhältnis konvergiert skalenweise
$2{,}734\to2{,}726\to2{,}721\to2{,}720\to2{,}719\to2{,}719$ gegen $e=2{,}71828$.

\section{Kontinuierliche, nicht diskrete Skaleninvarianz}
Die Selbstähnlichkeit ist \emph{kontinuierlich} -- eine log-Spirale, ein Potenzgesetz --, nicht
diskret log-periodisch. Ein Test auf diskrete Skaleninvarianz (das Residuum von $D(k)$ gegen das
glatte $\ln(1+k/75)$) liefert eine Amplitude nahe null und keine Periodizität in $\ln k$. Das
deckt sich mit dem Corpus: Dok.~146 hält fest, dass die Skalierung über den Faktor
$1/\xi_0=7500$ durch etwa $13$ logarithmische Sub-Level quasi-kontinuierlich wird und die
fraktale Dimension ohne Stufen konstant bleibt. Die Zeit-Windung realisiert dieselbe
kontinuierliche Skaleninvarianz geometrisch, als log-Spirale.

\section{Raum und Zeit: dasselbe $\xi_0$, verschieden skaliert}
Raum und Zeit tragen dieselbe Fundamentalkonstante $\xi_0$ an zwei Stellen: das räumliche
Box-Dimensions-Defizit ist $\xi_0$ (in $D_f=3-\xi_0$), der Windungs-Defekt je Gatter ist
$100\,\xi_0=1/75$. Der Faktor $100$ dazwischen ist der Rekursions-Koeffizient; er verbindet ein
\emph{geometrisches} Defizit mit einem \emph{dynamischen} Defekt und zeichnet keine Achse aus.
Welche Achse den Defekt trägt, ist Projektionssache -- er wird herkömmlich auf der Zeit notiert,
ist unter $\tilde T\cdot m=1$ aber symmetrisch (folgende Abschnitte), also keine Verstärkung der
Zeit gegenüber dem Raum.

Eine Grenze ist dabei ausdrücklich zu ziehen: die Box-Dimension der Zeit-Windung selbst ist
$\approx 1$ -- sie ist eine Kurve -- und \emph{nicht} $3-\xi_0$. Die Verbindung zwischen der
temporalen Spirale und der räumlichen fraktalen Dimension ist das geteilte $\xi_0$, nicht eine
Gleichheit der Dimensionen. Die Zeit-Spirale trägt nicht die Raumdimension.

\section{Die Achsenwahl ist eine Projektion}
Die logarithmische Spirale ist eine Eigenschaft der gesamten gekoppelten Windung auf $T^4$, nicht
einer einzelnen Koordinate. Erzwungen ist die \emph{Nichtschließung} mit Verhältnis $e$; auf
welche Achse man sie legt, ist Darstellungssache. Weil $\tilde T\cdot m=1$ Zeit und Masse invers
verkoppelt und der Torus alle Kreise verwebt -- nichts faktorisiert --, gibt es keine
ausgezeichnete Zeit-Dimension, die für sich allein aufrollt. Eine generische Windung läuft
gleichzeitig durch mehrere Kreise; die reine \emph{Zeit}-Windung ist der Spezialfall einer
Projektion. Dieselbe Nichtschließung ließe sich ebenso auf eine Raumachse oder auf eine Mischung
legen -- andere Basis, gleiche Invariante.

Die Zuordnung auf die Zeit ist damit eine künstliche Verschiebung auf eine Achse -- eine
Konvention, keine Auszeichnung. Der Gatter-Defekt wird herkömmlich auf der Zeit-Windung
beschrieben (Dok.~034); die duale Projektion (nächster Abschnitt) zeigt jedoch, dass dieselbe
Nichtschließung mit demselben Faktor $100$ auf der Raumseite erscheint. Keine Basis ist damit
bevorzugt; jede ist eine gleichwertige Darstellung derselben Invariante -- dieselbe
Darstellungsfreiheit wie zwischen den zwei äquivalenten Repräsentationen (Dok.~206) und in der
Translations-Invarianz von $\theta=2/9$ (Dok.~293).

\section{Kontinuierliche Zeit: die duale Projektion in den Raum}
Setzt man die Zeit umgekehrt \emph{kontinuierlich} an -- glatt, ohne gatterweises Aufrollen --,
so kann die Nichtschließung nicht mehr auf der Zeit sitzen; sie muss von einer anderen Achse
getragen werden. Die Bindung $\tilde T\cdot m=1$ lässt die Zeit nicht isoliert glätten: mit
$d\tilde T/\tilde T=-\,dm/m$ zieht die Masse invers mit, und mit ihr die räumlich-geometrische
Seite ($D_f=3-\xi_0$). Die laufende Windung erscheint dann als \emph{Variation des Raums} statt
als aufgerollte Zeit. Es ist dieselbe invariante Nichtschließung mit Verhältnis $e$, nur in die
duale Projektion gedreht: in der jetzigen Lesart läuft die Zeit (rollt auf) und der Raum bleibt
vergleichbar; in der dualen Lesart ist die Zeit kontinuierlich und der Raum variiert.

Diese Spiegelung ist durchgerechnet (\texttt{dual\_projektion\_probe.py}). Weil die Zeit je
Gatter um den Faktor $K_{\text{frak}}=1-100\,\xi_n$ verkürzt vorrückt, rückt die Masse nach
$\tilde T\cdot m=1$ um den inversen Faktor $1/K_{\text{frak}}$ vor; der Defekt je Schritt ist auf
beiden Seiten $100\,\xi_n$ (das Verhältnis $d_{\text{Masse}}/d_{\text{Zeit}}\to1$ auf sechs
Stellen). Der Faktor $100$ \emph{wandert} also nicht, er ist \emph{gespiegelt}: die
Gatter-Faktoren sind invers ($K_{\text{frak}}\leftrightarrow1/K_{\text{frak}}$), die Verstärkung
$100$ ist dieselbe. Beide Projektionen ergeben denselben logarithmischen Verlauf, und der
Spiral-Fit liefert für beide dasselbe Verhältnis $e$ ($e^{2\pi a}=2{,}7195$ bzw.\ $2{,}7192$);
ihre kumulierten Defekte unterscheiden sich nur um eine \emph{beschränkte} Konstante
($\approx0{,}0134$, wächst nicht mit $\ln k$). Der Kerngehalt ist damit darstellungsinvariant und
die Nichtschließung erhalten: man kann die Zeit nicht glätten, ohne dass Masse und Raum die
Windung übernehmen -- das Masse-Zeit-Dualitätsprinzip in Aktion. Ihre Verteilung auf Zeit, Masse
und Raum ist Projektionssache; das Verhältnis $e$ ist es nicht.

\section{Einordnung zu HLV}
Die FFGFT-interne Aussage ist damit \emph{abgeleitet}: die K$_{\text{frak}}$-Rekursion erzwingt die
Zeit-Windung als log-Spirale mit Verhältnis $e$, gleichwertig mit einem $1/n$-Zerfall ihrer
Ganghöhe. Erst hier tritt HLV hinzu -- und dabei ist vorab zu trennen, dass \enquote{Spiral-Zeit}
in HLV nicht ein Objekt bezeichnet, sondern zwei. Erstens die \emph{geometrische} Zeit aus der
Faktorisierung $T^7=T^4\times T^3$ (Dok.~285): dort ist die Zeit ein $A_5$-Singlet auf dem
$S^1$-Faktor, entkoppelt und separabel. Zweitens die \emph{temporal-dynamische} Spiral-Zeit, das
Objekt, das HLV so \emph{nennt}: ein Nicht-Markov-Operator mit triadischer Zeitkoordinate und
Gedächtniskanal, ohne $A_5$ und ohne $S^1$-Faktor. FFGFT faktorisiert im Gegensatz zu beiden nicht;
seine Spirale ist die Projektion eines untrennbaren Ganzen.

Der oben genannte Grenzfall konstanter Ganghöhe ($d=0$) entspricht dem Fixpunkt $\xi\to0$ der
Rekursion, in dem der Defekt verschwindet und die Windung schließt. Das ist die \emph{geometrische}
entkoppelte Singlet-Zeit, nicht das temporal-dynamische Objekt; die laufende FFGFT-Windung nähert
sich ihr asymptotisch. Die entscheidende Frage an HLV lautet damit nicht mehr \enquote{ähnlich?},
sondern \enquote{Verhältnis $e$, ja oder nein?} -- und sie ist an das richtige der beiden Objekte
zu richten. Der folgende Abschnitt führt diese Prüfung durch. In den drei methodischen Schichten
geordnet: die Ableitung Rekursion~$\to$~log-Spirale ist Kern-Ableitung, aus $\xi$ gewonnen; die
Gleichsetzung mit einer der beiden HLV-Zeiten ist ein Brücken-Kandidat; ein physikalischer Anspruch
über HLV wird nicht erhoben.

\section{Abgleich mit HLVs Spiral-Zeit: geteilte $Z_3$, kein Verhältnis $e$}
Die Gleichsetzung der FFGFT-Log-Spirale mit HLVs \enquote{Spiral-Zeit} lässt sich nun an einer
konkreten HLV-Quelle prüfen, und zwar an dem im vorigen Abschnitt als \emph{temporal-dynamisch}
bezeichneten der beiden Objekte. Formal ist es ein Nicht-Markov-Operator mit triadischer
Zeitkoordinate
\begin{equation}
  \Psi(t)=(t,\phi(t),\chi(t)),\qquad \Psi(t)=t+i\,\phi(t)+j\,\chi(t),
\end{equation}
wobei $t$ die Laborordnung, $\phi$ die Phasenorganisation und $\chi$ ein Gedächtnis notiert; es
trägt weder $A_5$ noch Singlet noch $S^1$-Faktor und besitzt einen Gedächtniskern $K_\chi$ mit
Markov-Grenzfall $g_\chi\to0$. Genau dieses Objekt nennt HLV \emph{Spiral-Time}; der entkoppelte
$A_5$-Singlet (vorige Abschnitte) ist das andere. Beide teilen nur den Namen.

\section{Das Verhältnis $e$ trägt das temporal-dynamische Objekt nicht}
Die oben formulierte Testfrage -- \enquote{Verhältnis $e$, ja oder nein?} -- zielt auf ein
geometrisch selbstähnliches Objekt. Das temporal-dynamische Objekt ist von dieser Kategorie nicht.
Es ist vielmehr ein Objekt der offenen Quantendynamik: seine Gesamtdynamik ist -- nach Marcels
Spiral-Time-Papier -- eine unitäre Dilatation auf dem Produktraum $H_S\otimes H_M$ aus System und
Gedächtnis, und geprüft wird es dort über Kriterien offener Systeme (Prozesstensor,
CP-Divisibilität, Hard-Reset, Finite-Bath-Vergleich), nicht über eine geometrische
Selbstähnlichkeit. Es gehört damit erkennbar der dynamischen, nicht der geometrischen Kategorie an;
die $e$-Frage ist an ein Objekt gerichtet, das sie nach eigener Konstruktion nicht beansprucht.
Es trägt kein abgeleitetes Selbstähnlichkeits-Verhältnis, keinen $1/n$-Zerfall einer Ganghöhe und
keinen Skalenlauf; seine Nichtschließung sitzt im Gedächtniskanal $\chi$, nicht in einer laufenden
Skala. Die einzige selbstähnliche interne Struktur ist die golden-ratio-/phasonische Ordnung des
zugehörigen Cut-and-Project (die Kontrolle ersetzt die $\varphi$-Struktur durch zufällige
irrationale Phasen), also $\varphi$ und nicht $e$. Auf dem $e$-Kriterium ist die Gleichsetzung
damit nicht gestützt. Was beide Rahmen strukturell teilen, ist die Rolle der Nichtschließung als
temporales Kernobjekt und ein expliziter Gedächtnis-aus-Grenzfall: FFGFTs $d=0$-Schließung
($\xi\to0$) entspricht Marcels Markov-Rückfall $g_\chi\to0$. Gleiche Rolle, verschiedener
Mechanismus -- geometrische Akkumulation ($D(k)=\sum 100\,\xi_n\to\ln$) gegen stochastischen
Gedächtniskern.

FFGFTs Nichtschließung liegt dabei nicht nur geometrisch als Log-Spirale vor, sondern gleichwertig
auch im Hilbertraum: als Gedächtniskern, die Fourier-Transformierte der diskreten Spektraldichte
$\sum_k g_k^2\,\delta(\omega-\omega_k)$ des $T^4$-Connection-Laplace (Dok.~283) -- dasselbe
untrennbare Objekt in zwei Darstellungen ($H=L^2(T^4)\otimes\mathbb{C}^3$, verlustfreie Bijektion,
Dok.~230/282); das ist Darstellungs-Robustheit, keine zusätzliche Evidenz. Damit lässt sich der
Abgleich in Marcels eigener Sprache führen, Kern gegen Kern: FFGFTs Kern ist oszillierend und
diskret-spektral (Revivals, BLP-Rückfluss $5{,}125$), Marcels abklingend und beschränkt --
dieselbe Objektklasse, verschiedene Gestalt (Dok.~283). Genau der $1/n$-Skalenlauf, der $e$ trägt,
fehlt dem abklingenden Kern.

\section{Rechnerische Prüfung des $e$-Kriteriums}
Das $e$-Kriterium ist notwendig und hinreichend über den Ganghöhen-Zerfall: eine Windung ist genau
dann äquiangular (log-Spirale, Verhältnis $e$ je $2\pi$), wenn der Defekt je Schritt wie $1/n$
zerfällt und der kumulierte Defekt logarithmisch wächst. Ein konstanter Defekt ergibt dagegen eine
Archimedische Windung mit Verhältnis $\to1$. Beide Konstruktionen lassen sich auf ihren eigenen
Definitionen daran messen (\texttt{spiralzeit\_verhaeltnis\_probe.py}).

FFGFT trägt $e$: mit $d_n=100\,\xi_n$ gilt $d_n\,(n+75)\to1$ -- also der $1/n$-Zerfall --, und der
Spiral-Fit liefert $e^{2\pi a}=2{,}71947\approx e$. Marcels Spiral-Zeit trägt es nicht: ihre Phase
$\phi_i=q_\phi\,u_i$ ist linear in der internen Koordinate und ihr Gedächtnis $\chi_i=F(\rho_i)$
ist ein beschränktes Fenster; der Defekt je Schritt bleibt konstant, $D(k)$ wächst linear statt
logarithmisch ($D(10^5)=1333$ gegen $\ln(1+k/75)=7{,}2$), und das gemessene Verhältnis ist
$1{,}0011$. Auf ihren eigenen Definitionen ist sie also gar keine log-Spirale, sondern eine Windung
konstanter Ganghöhe.

Die Geometrie, die HLV tatsächlich trägt, weist zudem auf eine andere Konstante als $e$. Baute man
aus dem $\tau=\varphi$-Cut-and-Project eine äquiangulare Spirale, so wäre ihr natürlicher
Selbstähnlichkeits-Faktor golden -- $\varphi=1{,}618$ je Umlauf bzw.\ $\varphi^4=6{,}854$ je $2\pi$
--, nicht $e=2{,}718$. Beide Rahmen kennen damit, soweit sie überhaupt Selbstähnlichkeit tragen,
verschiedene Konstanten: $e$ auf der FFGFT-Seite (aus dem $1/n$-Zerfall der $\xi$-Rekursion),
$\varphi$ auf der geometrischen HLV-Seite (aus dem Cut-and-Project). Als Brücken-Bedingung im Sinne
von P35 heißt das präzise: HLVs Spiral-Zeit trägt $e$ nur, wenn ihr Phasengesetz einen
$1/n$-Skalenlauf -- eine laufende innere Skala -- enthält; die vorliegenden Definitionen enthalten
ihn nicht. Das ist ein Kandidat, keine Ableitung.

\section{Basis und Dimensionen: gleiche Zahlen, verschiedene Bedeutung}
Ob beide Rahmen \emph{dieselbe} Basis tragen, lässt sich dimensionsweise ordnen. Die Zahlen
stimmen überein, die Bedeutung der Dimensionen nicht. Das \emph{Dreifache}: FFGFTs Drei ist das
interne $\mathbb{C}^3$ bzw.\ die $Z_3$-Ordnung auf $T^4/Z_3$ (Dok.~282) -- ein
Moden- und Generationsindex über dem Torus. Marcels Drei ist die triadische \emph{Zeit}
$(t,\phi,\chi)$ -- eine Zerlegung der einen Ordnungsachse in Ordnung, Phase und Gedächtnis. In
seiner Wurzelmassen-Notiz fällt diese Drei mit einem $C_3/Z_3$-Orbit ($2\pi k/3$, $k=0,1,2$)
zusammen; insofern teilen beide Rahmen die $Z_3$ ($C_3<A_5$, Dok.~285). Geteilt ist die $Z_3$,
nicht ein gemeinsamer dreidimensionaler Koordinatenraum: der eine indiziert Moden, der andere
Zeit-Rollen.

Das \emph{Sechsfache}: FFGFTs Sechs ist nicht nativ. Sie erscheint erst über die HLV-Einbettung
$T^7=T^4\times T^3$ (Dok.~285); FFGFT selbst ist vierdimensional ($T^4$). Marcels Sechs ist das
Elterngitter $\mathbb{Z}^6$ des Cut-and-Project, zerlegt in drei physikalische Richtungen
($P_\parallel$) und drei interne, phasonische ($P_\perp$). Dieselbe Zahl, verschiedene Anatomie:
eine $4+3$-Torus-Einbettung gegen ein $3+3$-Cut-and-Project. Zudem ist die Sechs bei Marcel nicht
die Basis der \emph{kontinuierlichen} Spiral-Zeit -- diese ist triadisch (die Drei) --, sondern
die des \emph{diskreten} $G$-Gitters; die beiden Ebenen sind bei ihm ausdrücklich Repräsentationen
desselben Operators, nicht dessen Ersatz. Das entspricht der Darstellungsfreiheit auf der
FFGFT-Seite (zwei äquivalente Repräsentationen, Dok.~206; aufgerollte $75$-Windung gegen entrollte
Log-Spirale).

Der einzige konkret belegte Überlapp ist das $\varphi$-ikosaedrische geteilte Objekt
($C_3<A_5$), rechnerisch etabliert, aber \emph{einseitig} -- die Prüfrichtung ist
FFGFT$\to$HLV (Dok.~294) -- und filtrations-empfindlich: es trennt HLV scharf von kristallinen
Kontrollen, aber nur schwach von generischem $\sqrt2$-Cut-and-Project. Rechnerisch lässt sich das am
Achsenspektrum festmachen (\texttt{basis\_6d\_achsen\_probe.py}): die sechs projizierten Achsen des
$6$D-Cut-and-Project sind nur bei $\tau=\varphi$ äquiangular -- alle fünfzehn paarweisen Winkel
gleich $63{,}435^\circ$ --, während $\tau=\sqrt2$ ein zweiwertiges Spektrum
($61{,}87^\circ/70{,}53^\circ$) ergibt. Die Sechs decken sich also nicht per se; gleich ist allein
die golden-ikosaedrische Projektion. Damit ist die Basis-Gleichheit ein \emph{Kandidat}, keine
abgeleitete Identität. Die Zahlen $3$ und $6$ decken sich, die geteilte Invariante ist die $Z_3$,
und die Dimensionen selbst tragen verschiedene Bedeutungen.

\section{Status}
Gesichert und abgeleitet ist: die Zeit-Windung als logarithmische Spirale mit
Selbstähnlichkeits-Verhältnis $e$, kontinuierliche Skaleninvarianz, aus der
K$_{\text{frak}}$-Rekursion; die exakte $75$-Schließung im eingefrorenen Fall; die geteilte
Konstante $\xi_0$, die als räumliches Box-Defizit $\xi_0$ und als Windungs-Defekt $100\,\xi_0$
erscheint (Faktor $100$ = Rekursions-Koeffizient, achsensymmetrisch). Die geschlossene $75$-Windung
ist dabei die aufgerollte (kompakte, Ein-Skalen-)Lesart, die Log-Spirale die entrollte
(dekompaktifizierte, Über-die-Skalen-)Lesart derselben Zeit-Windung.

Der Abgleich mit HLVs Spiral-Zeit ist auf dem $e$-Kriterium negativ: das temporal-dynamische
Objekt trägt kein Verhältnis $e$, sondern golden-ratio-/phasonische Quasiperiodizität; die
Basis-Frage reduziert sich auf die geteilte $Z_3$ ($C_3<A_5$, einseitig verifiziert,
filtrations-empfindlich). Die scharfe Restfrage verschiebt sich auf das \emph{geometrische}
HLV-Objekt -- ob die faktorisierte Schicht ($A_5$-Singlet auf $S^1$) ein skaleninvariantes Objekt
kennt, das gegen $e$ prüfbar wäre --; in den vorliegenden Quellen ist ein solches nicht
ausgewiesen. Die Kategorien-Trennung (zwei Spiral-Zeiten) und der dimensionsweise Abgleich sind
gesicherte Feststellungen; die Gleichsetzung bleibt ein Brücken-Kandidat, auf dem $e$-Kriterium
derzeit blockiert.

Der Kontrast ist dabei strukturell, nicht ontologisch. Beide Rahmen -- und ebenso die beiden
Varianten innerhalb von FFGFT, ob Zeit oder Raum aufgerollt wird -- sind mathematische Modelle,
nicht die Realität; keine Faktorisierung und keine Achsenwahl ist dadurch ontologisch
ausgezeichnet. Dass man hofft, ein solches Modell sei auf die Realität anwendbar, ist die
Motivation, keine Aussage darüber, was ist.

Reproduzierbar in \texttt{zeit\_windung\_probe.py}, \texttt{zeit\_logspirale\_probe.py},
\texttt{dual\_projektion\_probe.py}, \texttt{spiralzeit\_verhaeltnis\_probe.py} und
\texttt{basis\_6d\_achsen\_probe.py}.
